\section{CAMP Background and Relationship to WORKS}
\label{sec:design}
CAMP's measurement, prediction, and selective-feedback architecture was
introduced in WORKS~\cite{camp2026works}. We summarize the existing design
and its implementation to identify the components examined in this study.

\subsection{Measurement and compatible history}
\label{sec:history}
The eBPF tracer is attached to a task launcher and tracks descendants
through process-fork events. Kernel handlers discard events from
untracked execution contexts, preventing unrelated tasks on the same node
from contributing to that task's measurements. For explicit reads, the
implementation associates descriptors with observed file opens and
accumulates positive return values from \texttt{read} and \texttt{pread64}.
These are returned bytes, rather than requested buffer lengths; descriptor
attribution depends on the file operations observed by the tracer.

A successful \texttt{mmap} call establishes a virtual mapping but does not
show how much mapped data is accessed. CAMP instead instruments the
\texttt{filemap\_fault} entry point for file-backed page-fault activity.
The retained tracer adds 4,096 bytes per attributed event. The resulting
consumption signal is explicit-read bytes plus this page-fault byte
estimate. It does not deduplicate pages or certify that each entry event
completed a successful fault; it is therefore not an exact count of unique
bytes consumed, physical disk traffic, or resident memory. This distinction
motivates the read-only versus read-plus-mmap comparison. A task that
streams data and one that retains an index may have similar consumption
but different memory peaks.

Completed observations are organized into compatibility scopes.
$I_0$ and $I_1$ match task-instance/split-group and input identity,
respectively; $I_2$ selects nearby static-feature signatures; $I_3$ pools
process history. These scopes preserve workflow, version, and system
context. $I_4$ broadens to workflow history and $I_5$ to system-compatible
history. Support thresholds select a usable scope, while summaries from
multiple scopes remain features. History scopes identify relevant
observations; prediction-model selection is a separate mechanism.

\subsection{Consumption prediction and memory allocation}
\label{sec:consumption}
The Access baseline, denoted A+P in the artifact, uses static features and
peak-memory history. CAMP adds historical consumption and a prelaunch
consumption estimate:
\begin{align}
\mathbf{x}^{\mathrm{AP}}_t &= [\mathbf{s}_t,\mathbf{h}_{m,t}],\\
\hat c_t &= g(\mathbf{s}_t,\mathbf{h}_{m,t}),\label{eq:proxy}\\
\mathbf{x}^{\mathrm{CAMP}}_t
  &= [\mathbf{s}_t,\mathbf{h}_{m,t},\mathbf{h}_{c,t},\hat c_t].
\end{align}
The estimator $g$ learns from audited outcomes. Historical consumption
$\mathbf{h}_{c,t}$ describes earlier executions, whereas $\hat c_t$ estimates
the pending task's consumption. The four feature variants retain A+P
alone, add either signal separately, or include both. In the offline
training pipeline, cross-fitted proxy predictions supply the training
feature: each held-out group is predicted by a model fitted without that
group. Measured current-task consumption is not substituted for this
prediction.

\label{sec:quantiles}
LightGBM ensembles estimate peak-memory quantiles at
$\{0.50,0.90,0.95,0.99,0.995\}$. The process-tail configuration combines a
global central estimate with workflow--process upper-tail models and a
global fallback; global-tail configurations also exist in the artifact.
Feature comparisons must therefore hold the prediction configuration
fixed or explicitly identify configuration differences. Changing features
and model selection together cannot isolate the contribution of a feature.

\subsection{Calibration and selective feedback}
\label{sec:calibration}
A development-selected policy chooses a base estimate $b_t$ and a residual
quantile. Hierarchical calibration uses log residuals
$r_\tau=\log(y_\tau/b_\tau)$ and blends a supported residual scope with a
broader parent. Residual support determines this scope independently of
the history-feature scope. The correction $\rho_t$ produces
\begin{equation}
\hat y_t=\left\lceil b_t\exp\!\left(\max(0,\rho_t)\right)\right\rceil_{\mathrm{MiB}}.
\label{eq:final}
\end{equation}
The adjustment is one-sided and rounded to whole mebibytes. Cross-fitted
predictions provide development residuals; subsequently available outcomes
support online updates. Calibration remains empirical and provides no
unconditional coverage guarantee on shifted inputs.

\label{sec:gates}
The selective policy allocates an audit budget using prelaunch risk,
information, and discovery signals. Peak-memory outcomes can update
history after completion; consumption outcomes enter only when measured.
The deployment mechanisms are distinct from this policy: the generic
Nextflow audit overlay supports process-name filtering and periodic task
selection, while the retained live Bowtie2 configuration audits every task.
Thus, the live run does not by itself demonstrate selective-auditing gains.

In the artifact, Nextflow configuration overlays add instrumentation to
unmodified nf-core workflow sources. A launch wrapper pauses the target,
attaches the tracer, waits for readiness, and resumes execution. Bowtie2's
allocation integration additionally defines process-specific memory
closures that pass task metadata to the prediction service and return
requests to Nextflow for Slurm submission. Completion hooks record
outcomes, and an asynchronous learner publishes updated model versions
for later predictions. The retained launcher requests model updates after
each 128 eligible outcomes; history ingestion and model retraining are
separate operations. These additions preserve analytical tool commands
but require workflow-specific metadata bindings, deployment configuration,
and Linux hosts with BCC/eBPF support and tracing privileges. They are not
a transparent allocator for every workflow engine.

\subsection{Prior evidence and extension boundary}
WORKS already reports the architecture, Access and Sizey comparisons,
selected feature and mmap ablations, and selective-auditing
results~\cite{camp2026works}. Earlier Bowtie2 records provide repeated-input deployment evidence;
the new executions are packaged separately in the artifact~\cite{campArtifact}. The present study develops the
component analysis, examines downstream allocation benefit relative to
audit cost, and
evaluates online adaptation to new inputs and the effect of an additional
warm pass from a shared cold checkpoint. Reanalysis of existing
measurements is distinguished from new experiments. Routing, calibration,
and compatibility scopes are described as background; separate ablation
sweeps of these mechanisms are outside the present scope.
